%% precise-ledger-pro-arxiv.tex
%% arXiv submission: cs.DB (primary), cs.SE (secondary)
%% Authors: Vibhor Mahajan, Raman Bedi
%% Draft — April 1, 2026
%%
%% REVIEW FIXES APPLIED (from peer review, April 2026):
%%   D1  – Replaced fabricated Meijer & Fokkinga [1991] citation with
%%          Gray & Reuter (1992) "Transaction Processing".
%%   D2  – Replaced misattributed Hellerstein et al. [2019] with
%%          Gray & Reuter (same citation, legitimately relevant).
%%   D3  – Removed dangling Graefe [2011] reference; added in-text
%%          citation in §8.4 where it is substantively relevant.
%%   D4  – Abstract revised: "no balance drift" → "no unintended drift".
%%          Added clarifying sentence in §7.3.
%%   D5  – Raman Bedi affiliation placeholder added (confirm before
%%          submission).
%%   D6  – Algorithm 1 complexity: distinct variable p introduced for
%%          fork-prefix size to resolve k/k name collision.
%%   D7  – Draft metadata moved entirely to these comment lines; absent
%%          from rendered PDF.
%%   D8  – ACM CCS primary corrected to "Transaction management".
%%   D9  – Trailing space and bare URL fixed; \url{} used.
%%   D10 – Single-entry model limitation stated in §8.4.
%%   D11 – Commercial platform names removed; replaced with generic text.
%%   D12 – HTML <sub> tags replaced with proper LaTeX subscripts.

\documentclass[11pt,letterpaper]{article}

\usepackage[T1]{fontenc}
\usepackage[utf8]{inputenc}
\usepackage{times}                  % arXiv-recommended: ~10% more text per page
\usepackage[margin=1in]{geometry}
\usepackage{amsmath,amssymb}
\usepackage{booktabs}               % professional table rules
\usepackage{tabularx}
\usepackage{array}
\usepackage{microtype}              % better line-breaking
\usepackage{parskip}                % inter-paragraph spacing
\usepackage[hyphens]{url}
\usepackage[colorlinks=true,
            linkcolor=blue,
            citecolor=blue,
            urlcolor=blue]{hyperref}
\usepackage{algorithm}
\usepackage{algpseudocode}
\usepackage{listings}
\usepackage{xcolor}
\usepackage{enumitem}
\usepackage{caption}
%% subcaption removed (R4-4): imported but never used

%% Listings style for pseudocode-like verbatim blocks
\lstset{
  basicstyle=\ttfamily\small,
  breaklines=true,
  columns=flexible,
  keepspaces=true,
  frame=single,
  framesep=4pt,
  rulecolor=\color{gray!40},
  backgroundcolor=\color{gray!5},
}

%% Note: the built-in \Comment{} from algpseudocode is used throughout.

\title{\textbf{Retroactive Reconciliation in Financial Ledgers:\\
A Fork-and-Sync Architecture for Out-of-Order Activities\\
and Reversals with Per-Entry Traced Corrections}}

\author{
  Vibhor Mahajan\\
  \small Independent Researcher\\
  \small \texttt{vibhor.mahajan@gmail.com}
  \and
  Raman Bedi\\
  \small \textit{[Affiliation — confirm before submission]}\\
}

\date{April 1, 2026}

\begin{document}

\maketitle

\begin{abstract}
Financial ledger systems in lending institutions must accommodate two
classes of late-arriving events that violate chronological ordering
assumptions: out-of-order (backdated) payment activities and reversals
of previously applied activities. Existing approaches---full ledger reconstruction, append-only
compensating entries, and mutable rewriting---each fail to
simultaneously preserve the complete audit trail, individually
attribute each correction entry to its source activity, maintain
derived-calculation correctness, and support arbitrary future
retroactive events. We present a \emph{fork-and-sync} architecture
that addresses all four requirements. When a retroactive
event arrives, the system forks the primary ledger at the appropriate
historical insertion point, applies the event and replays all subsequent
activities on the fork, and then reconciles the fork back into the
primary ledger by computing per-activity balance-impact deltas. Each
reconciliation emits individual, source-attributed adjustment entries
rather than a single grouped correction, preserving entry-level audit
traceability and enabling clean composition of future retroactive
events. A virtual ledger clock decouples logical time from wall-clock
time to guarantee deterministic replay. We describe the algorithmic
design, implement it in a reference system, and provide a comparative
evaluation against a conventional grouped-adjustment ledger processing
nine activities that include a reversal and a backdated payment,
demonstrating full audit traceability with no unintended balance drift.
\end{abstract}

%% ACM CCS 2012 (D8 fix)
\noindent\textbf{ACM CCS:} Information systems~$\to$~Transaction management;
Information systems~$\to$~Data streams;
Software and its engineering~$\to$~Software reliability.

\noindent\textbf{Keywords:} financial ledger, retroactive reconciliation,
fork-and-sync, audit trail, idempotency, bitemporal, event replay.

\tableofcontents
\newpage

% =========================================================================
\section{Introduction}
% =========================================================================

Loan servicing systems must process millions of ledger activities daily,
yet financial events routinely arrive out of chronological order. A
payment may be applied days after its contractual effective date, a
disbursement may be discovered to have been incorrectly processed, or a
regulatory investigation may require reversing a fee applied months ago.
Each such event is not merely an append at the tail of the ledger---it
changes the basis on which every subsequent derived value (daily
interest accrual, balance allocation) was computed.

Three strategies are prevalent in practice. \textbf{Full
reconstruction} rebuilds the entire ledger from its initial state each
time a retroactive event arrives. This is correct but computationally
expensive: for a ledger with thousands of entries and a retroactive
event at position~$k$, the cost is proportional to the number of
entries following~$k$ that must be recomputed and persisted.
\textbf{Append-only compensation} appends a single correcting entry (or
a small group of entries) without recomputing the downstream derived
values that depended on the corrected balance. This is efficient but
produces balance drift whenever the corrected activity altered the basis
for any subsequent temporal calculation. \textbf{Mutable rewriting}
updates historical entry values in place, which destroys the immutable
audit trail required by financial regulations and introduces concurrency
hazards.

A fourth technique---grouping downstream corrections into a single
opaque adjustment entry---is balance-correct but violates a subtler
invariant: it destroys entry-level audit traceability. When a grouped
adjustment of, say, $+\$3.30$ appears in a ledger, neither an auditor
nor a subsequent retroactive event processor can determine which of the
original entries it corrects, or by how much, without replaying the
entire history from scratch.

This paper presents \textbf{PreciseLedgerPro}, a system and
accompanying algorithm family that resolves this tension. The core
innovation is a fork-and-sync protocol in which:
\begin{enumerate}[label=\arabic*.]
  \item The primary ledger is never mutated; its original entries form
        the permanent audit record.
  \item A retroactive fork contains the corrected view of history,
        computed efficiently from the insertion point onward.
  \item Per-activity differential reconciliation emits one individually
        source-attributed adjustment entry per divergent downstream
        activity---not a single grouped correction.
  \item A virtual ledger clock advances only forward, providing a
        consistent time reference for deterministic replay of
        time-dependent calculations.
  \item Composite idempotency keys prevent duplicate entries during
        recursive replay cycles.
\end{enumerate}

The result is a ledger in which every row---including every
correction---can be explained by reference to a single source activity,
and in which any subsequent retroactive event can be applied cleanly
without reference to a decomposed adjustment history.

The remainder of this paper is organized as follows. Section~\ref{sec:background}
provides background on financial ledger data models and formalizes the
problem. Section~\ref{sec:related} surveys related work.
Section~\ref{sec:architecture} describes the system architecture and
data model. Section~\ref{sec:algorithms} presents the four core
algorithms. Section~\ref{sec:implementation} describes the reference
implementation. Section~\ref{sec:evaluation} provides a comparative
evaluation. Section~\ref{sec:discussion} discusses design trade-offs,
limitations, and future work. Section~\ref{sec:conclusion} concludes.

% =========================================================================
\section{Background and Problem Statement}
\label{sec:background}
% =========================================================================

\subsection{Ledger Data Model}

A financial ledger for a loan account maintains an ordered sequence of
\emph{ledger entries}. Each entry is an immutable record capturing:

\begin{itemize}
  \item An incremental balance change across multiple components:
        principal~$\delta_P$, interest~$\delta_I$, fee~$\delta_F$, excess~$\delta_E$.
  \item A running balance $(P, I, F, E)$ after this entry, computed as
        the running balance of the preceding entry plus the incremental
        change.
  \item An \texttt{effectiveAt} timestamp: the business-effective date
        of the originating activity.
  \item A \texttt{createdAt} timestamp: the system processing time.
  \item A \emph{source activity identifier}: a (type, id) pair that
        links the entry to the ledger activity that created it.
\end{itemize}

A \textbf{ledger activity} is an abstract event---a payment, a
disbursement, a daily interest accrual, a reversal---that \emph{applies
to} a ledger by generating one or more ledger entries.

Formally, let $L = \langle e_1, e_2, \ldots, e_n \rangle$ be the
primary ledger ordered by \texttt{effectiveAt}. For a new activity~$a$
with effective date~$t$:
\begin{itemize}
  \item $a$ is \textbf{in-order} if $t \geq \text{effectiveAt}(e_n)$.
  \item $a$ is \textbf{backdated} (out-of-order) if
        $t < \text{effectiveAt}(e_n)$.
\end{itemize}

\subsection{Derived Values and the Correction Problem}

Daily interest accrual---the canonical temporal ledger activity---computes:
\begin{equation}
  \text{interest}_d = P_{d-1} \times \frac{r}{d_c}
\end{equation}
where $P_{d-1}$ is the principal balance at the close of day $d-1$,
$r$~is the annual interest rate, and $d_c$~is the day-count basis
(e.g., $d_c = 365$ for actual/365). If a
backdated event at position~$k$ changes $P_{k-1}$, then every
subsequent daily interest entry must be recomputed, since each entry's
interest accrual depends on the running principal balance at its own
computation time.

The correction problem is therefore: given a retroactive event that
changes the balance at position~$k$ in a ledger of $n$ entries
($k < n$), efficiently propagate the correct balance deltas to all
entries at positions $k{+}1$ through $n$ without destroying the
original entry history.

\subsection{Requirements}

We derive four requirements from financial system practice:

\begin{description}
  \item[R1 (Correctness)] After applying any retroactive event, running
    balances of all subsequent entries must reflect the corrected history.
  \item[R2 (Audit Immutability)] Original ledger entries must not be
    modified or deleted.
  \item[R3 (Entry-Level Traceability)] Each correction entry must
    identify the source activity it corrects.
  \item[R4 (Retroactive Composability)] The ledger must support a
    subsequent retroactive event without requiring decomposition of
    prior correction entries.
\end{description}

Systems based on full reconstruction satisfy R1--R4 but at high
computational cost. Grouped-adjustment approaches satisfy R1--R2 but
violate R3--R4. Mutable rewriting satisfies only R1. This paper
provides a system that satisfies all four.

% =========================================================================
\section{Related Work}
\label{sec:related}
% =========================================================================

\subsection{Event Sourcing and Append-Only Ledgers}

Event sourcing \cite{fowler2005} stores state changes as an immutable
sequence of events from which current state is derived by replay. While
event sourcing provides R2, standard implementations reconstruct current
state by full replay, satisfying R1 only if the replay is complete and
ordered. They do not address the efficiency problem of partial
re-derivation from an insertion point, nor do they produce individually
traced correction entries (R3).

\subsection{Bitemporal Data Models}

Snodgrass~\cite{snodgrass1999} and Jensen et al.~\cite{jensen1994}
formalize bitemporal databases that track both \emph{valid time}
(business-effective time) and \emph{transaction time} (system recording
time). Bitemporal models preserve full history but typically require
query-time reconstruction to derive correct values at arbitrary time
points---they do not produce an updated, materially correct running
ledger in place. The fork-and-sync approach presented here can be
viewed as a form of bitemporal reconciliation that materializes
corrected running balances in the primary ledger while respecting
transaction-time immutability.

\subsection{Compensating Transactions and Sagas}

The saga pattern \cite{garciamolinasalem1987} uses compensating
transactions to undo the effects of failed or rolled-back long-running
transactions. A reversal activity in our model is structurally similar
to a compensating transaction, but differs in a critical respect: a
simple compensation negates only the direct impact of the reversed
activity. It does not account for the downstream derived effects---the
interest accruals and allocation changes that occurred because of the
now-reversed balance state. Our reversal algorithm generates both a
compensation entry and a differential reconciliation of all downstream
derived values.

\subsection{Financial Transaction Processing}
\label{sec:related-financial}

Gray and Reuter~\cite{grayreuter1992} provide the foundational
treatment of transaction processing systems, including isolation,
durability, and compensating transactions. The OLAP literature
addresses aggregate balance queries over large financial datasets, but
handling retroactive insertions with maintained running balances has
received limited attention in academic systems research. Proprietary
commercial loan servicing platforms handle backdating operationally, but
their algorithms are not publicly documented.

\subsection{Idempotent Processing}

Idempotency in distributed systems is well-studied
\cite{bernsteingoodman1981}. Our use of composite idempotency keys on
ledger entries, rather than at the transaction level, is specific to the
recursive replay problem: the same ledger activity may be encountered
multiple times during nested retroactive recursion, and the key must be
stable across all encounters to prevent duplicates.

% =========================================================================
\section{System Architecture}
\label{sec:architecture}
% =========================================================================

\subsection{Core Data Structures}

\textbf{Ledger.}~An ordered list of \texttt{LedgerEntry} records
associated with a single loan account. The ledger supports:
\begin{itemize}
  \item \texttt{addEntry(e)}---idempotent append with running-balance
        recomputation.
  \item \texttt{rollbackToEntryBefore(effectiveAt)}---returns a new
        ledger (fork) containing only entries with
        \texttt{effectiveAt}~$\leq$~the given timestamp.
  \item \texttt{rollbackToEntryBefore(activityType, activityId)}---returns
        a fork excluding all entries from the first entry of a given
        source activity onward.
  \item \texttt{calculateTotalImpact(activityType, activityId)}---returns
        the sum of all balance changes across all entries attributed to
        a given source activity.
  \item \texttt{getCurrentBalance()}---returns the running
        multi-component balance.
\end{itemize}

The current model is a single-entry running-balance record per loan
account and does not enforce double-entry bookkeeping constraints
(debits = credits). Extension to double-entry contra-account ledgers is
future work.

\textbf{LedgerEntry.}~An immutable record; key fields are shown in
Table~\ref{tab:ledgerentry}.

\begin{table}[h!]
\centering
\caption{Key fields of \texttt{LedgerEntry}.}
\label{tab:ledgerentry}
\begin{tabularx}{\linewidth}{>{\raggedright\arraybackslash}p{5.2cm}X}
  \toprule
  \textbf{Field} & \textbf{Description} \\
  \midrule
  \texttt{entryType} & Semantic classification: Disbursal, Payment,
                       StartOfDay, Reversal, Adjustment \\
  \texttt{principal}, \texttt{interest}, \texttt{fee}, \texttt{excess}
                     & Incremental balance changes per component \\
  \texttt{principalBalance}, \texttt{interestBalance}, \ldots
                     & Running balances after this entry \\
  \texttt{effectiveAt} & Business-effective timestamp \\
  \texttt{createdAt}   & Processing-time timestamp \\
  \texttt{sourceLedgerActivityType} & Type of the originating activity \\
  \texttt{sourceLedgerActivityId}   & ID of the originating activity \\
  \bottomrule
\end{tabularx}
\end{table}

The composite idempotency key is a deterministic hash of
(\texttt{loanId}, \texttt{entryId}, \texttt{entryType},
\texttt{sourceLedgerActivityType}, \texttt{sourceLedgerActivityId},
\texttt{effectiveAt}).

\textbf{Balance.}~A 4-tuple $(P,\,I,\,F,\,E)$ supporting
component-wise addition, subtraction, negation, and
total-amount computation.

\subsection{Ledger Activities}

A \texttt{LedgerActivity} is an abstract object with
\texttt{effectiveAt}, \texttt{createdAt}, and \texttt{activityId}
fields. An activity \emph{applies to} a ledger by generating ledger
entries. Three concrete subtypes are defined:

\begin{itemize}
  \item \textbf{Transaction.}~A payment or disbursement that applies a
        \texttt{TransactionSpreadStrategy} to the current balance to
        determine per-component impacts. Two spread strategies are
        provided: \emph{computational spread} (configurable waterfall
        allocation with overflow to excess) and \emph{static spread}
        (direct per-component assignment).
  \item \textbf{TemporalActivity.}~A time-driven activity
        (e.g., daily interest accrual) that computes entries using an
        injected \texttt{TemporalActivityContext} supplying interest
        rates and day-count conventions. Injection enables deterministic
        replay independent of external state.
  \item \textbf{ReversalActivity.}~An activity that reverses a
        previously applied activity; handled by
        Algorithm~\ref{alg:reversal} (Section~\ref{sec:alg-reversal}).
\end{itemize}

An activity is classified as backdated if its \texttt{effectiveAt}
precedes the \texttt{effectiveAt} of the ledger's most recent entry.

\subsection{LedgerClock}

The \texttt{LedgerClock} is a virtual monotonic clock passed as a
parameter to all activity application methods. It:
\begin{itemize}
  \item Starts at a defined minimum timestamp.
  \item Advances only forward through \texttt{advanceTime(t)}.
  \item Rejects backward adjustments during forward processing.
\end{itemize}

During retroactive replay, the clock's \texttt{now()} value is fixed to
the snapshot time of the primary ledger at the point of invocation.
This ensures that time-dependent derived values (e.g., daily interest
amounts) are computed deterministically regardless of wall-clock time.
The clock's externally controlled nature also enables isolated,
reproducible unit testing of all processing paths.

\subsection{Component Interaction}

Figure~\ref{fig:orchestration} illustrates runtime component
interaction. The \texttt{LedgerService} orchestrates all algorithm
variants. Activity classification is performed on each iteration. For
normal activities, the activity is applied directly. For backdated
activities, \texttt{performBackdatedActivity} is invoked. For
reversals, \texttt{reverseLedgerActivity} is invoked. Both delegates
conclude with \texttt{syncWithRetroactiveLedger}.

\begin{figure}[h!]
\begin{lstlisting}
LedgerService
  applyLedgerActivities(ledger, activities, clock, ctx)
     |
     +-- [in-order]   --> activity.applyTo(ledger, clock, ctx)
     |
     +-- [backdated]  --> performBackdatedActivity(...)
     |                      |
     |                      +-- rollbackToEntryBefore(t)
     |                      +-- activity.applyTo(fork, ...)
     |                      +-- applyLedgerActivities(fork,...)  <- recursive
     |                      +-- syncWithRetroactiveLedger(...)
     |
     +-- [reversal]   --> reverseLedgerActivity(...)
                            |
                            +-- calculateTotalImpact(...)
                            +-- generateCompensationEntry(...)
                            +-- rollbackToEntryBefore(type, id)
                            +-- applyLedgerActivities(fork,...)  <- recursive
                            +-- syncWithRetroactiveLedger(...)
\end{lstlisting}
\caption{\texttt{LedgerService} orchestration and recursive delegation
to fork-and-sync sub-routines.}
\label{fig:orchestration}
\end{figure}

% =========================================================================
\section{Algorithms}
\label{sec:algorithms}
% =========================================================================

\subsection{Algorithm 1: Retroactive Fork-and-Sync for Backdated Activities}
\label{sec:alg-backdated}

\textbf{Precondition:} A ledger activity~$a$ arrives with
$a.\texttt{effectiveAt} < \text{ledger.latestEntry.effectiveAt}$.

\textbf{Postcondition:} The primary ledger contains the correct running
balances for all entries following the insertion point, with each
downstream correction attributed to its source activity.

\begin{algorithm}[H]
\caption{Retroactive Fork-and-Sync for a Backdated Activity}
\label{alg:backdated}
\begin{algorithmic}[1]
\Require Primary ledger $L$, backdated activity $a$, clock, context ctx
\Ensure $L$ updated with $a$ applied at its economic date and all
        downstream corrections individually traced
\State $\text{fork} \leftarrow L.\texttt{rollbackToEntryBefore}(a.\texttt{effectiveAt})$
\State $p \leftarrow |\text{fork.entries}|$
  \Comment{bookmark: fork-prefix size}
\State $a.\texttt{applyTo}(\text{fork},\,\text{clock},\,\text{ctx})$
\For{each new entry $e$ in $\text{fork}[p..\text{end}]$}
  \State $L.\texttt{addEntry}(e)$ \Comment{mirror direct entries to primary}
\EndFor
\State $R \leftarrow \texttt{getLedgerActivities}(L.\text{loanId},\;
       \texttt{effectiveAt} \geq a.\texttt{effectiveAt},\;
       \texttt{createdAt} \leq \text{clock.now})$
  \Comment{sorted by effectiveAt asc}
\State $\texttt{APPLY-LEDGER-ACTIVITIES}(\text{fork},\,R,\,\text{clock},\,\text{ctx})$
  \Comment{recursive}
\State $\texttt{SYNC}(L,\,\text{fork},\,\text{clock.now},\,
       a.\texttt{createdAt},\,p)$
\end{algorithmic}
\end{algorithm}

\paragraph{Note on recursion.}
Step~7 may itself encounter a backdated activity relative to
$\text{fork}$, triggering a nested invocation of
Algorithm~\ref{alg:backdated}. Idempotency keys applied at Step~5
(the \texttt{addEntry} call inside the For loop) ensure that entries
generated across recursion depths are not duplicated.

\paragraph{Complexity.}
Let $p$ be the fork-prefix size (number of preserved entries) and $m$
the number of replay activities. The fork copy is $O(p)$. Replay
processes $m$ activities, each $O(1)$ for direct entries; temporal
activities with waterfall spread are $O(c)$ where $c$ is the number of
balance components (constant in practice, $c = 4$). Reconciliation is
$O(m)$ with one pass per distinct source activity. Sorting the
replay activities at Step~6 dominates: total complexity per backdated
event is
$O(p + m \log m)$.

\subsection{Algorithm 2: Compensation + Retroactive Fork-and-Sync for Reversals}
\label{sec:alg-reversal}

\textbf{Precondition:} A reversal activity~$r$ identifies a target
activity~$t$ to be reversed. All entries attributed to~$t$ exist in
the primary ledger.

\textbf{Postcondition:} The primary ledger contains a compensation
entry immediately negating $t$'s impact, plus individually traced
adjustments for all downstream derived effects, with no record of
$t$'s existence removed from the ledger.

\begin{algorithm}[H]
\caption{Compensation + Retroactive Fork-and-Sync for a Reversal}
\label{alg:reversal}
\begin{algorithmic}[1]
\Require Primary ledger $L$, reversal activity $r$, clock, context ctx
\Ensure $L$ updated with compensation entry and per-source-activity
        downstream adjustment entries
\State $\Delta \leftarrow L.\texttt{calculateTotalImpact}(r.\text{targetType},\,r.\text{targetId})$
\State $\text{comp} \leftarrow \texttt{buildEntry}(\text{change} =
       {-}\Delta,\; \text{balance} = L.\text{currentBalance} - \Delta,\;
       \text{type} = \text{``Reversal''},\; \text{source} =
       (r.\text{targetType},\,r.\text{targetId}),\; \ldots)$
\State $L.\texttt{addEntry}(\text{comp})$
  \Comment{immediately corrects current balance}
\State $\text{fork} \leftarrow L.\texttt{rollbackToEntryBefore}(r.\text{targetType},\,r.\text{targetId})$
\State $p \leftarrow |\text{fork.entries}|$
\State $R \leftarrow \texttt{getLedgerActivities}(L.\text{loanId},\;
       \texttt{createdAt} > \text{originalActivity}(r.\text{targetId}).\texttt{createdAt},\;
       \texttt{createdAt} \leq \text{clock.now})$
  \Comment{excluding $r.\text{targetId}$; sorted by effectiveAt asc}
\State $\texttt{APPLY-LEDGER-ACTIVITIES}(\text{fork},\,R,\,\text{clock},\,\text{ctx})$
  \Comment{recursive}
\State $\texttt{SYNC}(L,\,\text{fork},\,\text{clock.now},\,
       r.\texttt{createdAt},\,p)$
\end{algorithmic}
\end{algorithm}

\paragraph{Exclusion invariant.}
The reversed target activity is excluded from replay at Step~6. This
prevents infinite recursion and ensures the fork represents the world
``as if the reversed activity never happened.''

\paragraph{Compensation semantics.}
The compensation entry at Step~3 immediately corrects the primary
ledger's current balance without waiting for reconciliation. This is
essential for balance queries that may occur between the compensation
step and the completion of replay.

\subsection{Algorithm 3: Differential Impact Reconciliation (SYNC)}
\label{sec:alg-sync}

This sub-routine is shared by Algorithms~\ref{alg:backdated}
and~\ref{alg:reversal}. It reconciles per-activity impacts between the
retroactive fork and the primary ledger, generating one adjustment entry
per divergent source activity and zero entries for activities whose
balance impact did not change.

\begin{algorithm}[H]
\caption{Differential Impact Reconciliation (SYNC)}
\label{alg:sync}
\begin{algorithmic}[1]
\Require Primary ledger $L$, retroactive fork $F$, current time $\tau$,
         adjustment effective date $t_{\text{adj}}$, fork index $p$
\Ensure $L$ updated with one adjustment entry per divergent source activity
\State $F_{\text{entries}} \leftarrow F.\text{entries sorted by effectiveAt}$
\State $L_{\text{entries}} \leftarrow L.\text{entries sorted by effectiveAt}$
\If{$F_{\text{entries}}$ is empty \textbf{or} $L_{\text{entries}}$ is empty}
  \State \Return
\EndIf
\State $\text{processed} \leftarrow \{\}$
\For{$i$ from $p$ to $|F_{\text{entries}}| - 1$}
  \State $e_F \leftarrow F_{\text{entries}}[i]$
  \State $\kappa \leftarrow (e_F.\texttt{sourceLedgerActivityType},\;
         e_F.\texttt{sourceLedgerActivityId})$
  \If{$\kappa \in \text{processed}$}
    \State \textbf{continue}
  \EndIf
  \State $\delta_L \leftarrow L.\texttt{calculateTotalImpact}(\kappa)$
  \If{$\delta_L = \text{null}$}
    \State $L.\texttt{addEntry}(e_F)$ \Comment{net-new activity}
  \Else
    \State $\delta_F \leftarrow F.\texttt{calculateTotalImpact}(\kappa)$
    \If{$\delta_F \neq \delta_L$}
      \State $\Delta \leftarrow \delta_F - \delta_L$
      \State $\text{adj} \leftarrow \texttt{buildAdjustmentEntry}(
             \text{change} = \Delta,\;
             \text{balance} = L.\text{currentBalance} + \Delta,\;
             \text{source} = \kappa,\;
             \texttt{effectiveAt} = t_{\text{adj}},\;
             \texttt{createdAt} = \tau)$
      \State $L.\texttt{addEntry}(\text{adj})$
    \EndIf
  \EndIf
  \State $\text{processed}.\text{add}(\kappa)$
\EndFor
\end{algorithmic}
\end{algorithm}

\paragraph{Key properties.}
\begin{itemize}
  \item Zero-delta activities produce no adjustment entry
        (``skip-if-equal'' guarantee).
  \item Each source activity is processed exactly once per
        reconciliation round (the \texttt{processed} set).
  \item Every adjustment entry carries full source-activity attribution,
        satisfying R3.
  \item The algorithm is $O(k)$ in the number of post-fork entries.
\end{itemize}

\subsection{Algorithm 4: Idempotent Entry Addition}
\label{sec:alg-idempotent}

\begin{algorithm}[H]
\caption{Idempotent Entry Addition}
\label{alg:addentry}
\begin{algorithmic}[1]
\Require Ledger $L$, entry $e$
\Ensure $L$ contains $e'$ (clone of $e$ with recomputed running balances); unchanged if $e$ was already present (duplicate)
\State $e' \leftarrow \texttt{clone}(e)$
\State $e'.\texttt{updateRunningBalances}(L.\texttt{currentBalance})$
\State $\kappa \leftarrow \texttt{idempotencyKey}(e')$
\If{any existing entry in $L$ has key $\kappa$}
  \State \Return \Comment{silent skip -- duplicate}
\EndIf
\State $L.\text{entries}.\texttt{append}(e')$
\end{algorithmic}
\end{algorithm}

The composite idempotency key is a deterministic function of
(\texttt{loanId}, \texttt{entryId}, \texttt{entryType},
\texttt{sourceLedgerActivityType}, \texttt{sourceLedgerActivityId},
\texttt{effectiveAt}). This key is stable across invocations: revisiting
the same activity during a nested retrospective cycle always produces
the same key, preventing duplication without relying on sequence numbers.

% =========================================================================
\section{Implementation}
\label{sec:implementation}
% =========================================================================

The reference implementation uses Java~21 with the Quarkus framework
(version 3.11) for the server-side service layer. Persistence uses
Hibernate/Panache with H2 (test/development) and PostgreSQL (production).
Monetary arithmetic uses the JavaMoney/Moneta library to avoid
floating-point precision loss on financial quantities; balance components
are stored as arbitrary-precision \texttt{NUMERIC} decimal values in the
database and manipulated via \texttt{MonetaryAmount} objects at runtime.

Key design decisions:

\begin{itemize}
  \item \textbf{Hexagonal architecture.}~The core algorithmic logic
        (\texttt{LedgerService}, \texttt{Ledger}, \texttt{LedgerActivity}
        hierarchy) contains no repository or framework dependencies.
        Repository interfaces are injected, enabling deterministic unit
        testing with in-memory ledger objects.
  \item \textbf{Injected temporal context.}~\texttt{TemporalActivityContext}
        supplies interest-rate schedules and day-count conventions as
        constructor parameters rather than runtime lookups. This is
        required for deterministic replay: a \texttt{StartOfDay} activity
        replayed on a retroactive fork must produce the same interest
        amount as it would have in forward processing.
  \item \textbf{Virtual clock injection.}~\texttt{LedgerClock} is passed
        as a parameter, not obtained from \texttt{System.currentTimeMillis()}.
        This decouples algorithmic correctness from wall-clock time and
        makes all processing paths directly unit-testable~\cite{feathers2004}.
  \item \textbf{Activity factory.}~A \texttt{LedgerActivityFactory}
        converts \texttt{GeneralLedgerActivity} persistence objects into
        typed \texttt{LedgerActivity} instances, isolating the
        persistence model from the domain model.
\end{itemize}

The implementation, including all test fixtures and integration tests,
is publicly available at
\url{https://github.com/v1bh0r/precise-ledger-pro}.

% =========================================================================
\section{Evaluation}
\label{sec:evaluation}
% =========================================================================

\subsection{Scenario Design}

We evaluate both the correctness and traceability properties of the
fork-and-sync approach against a conventional grouped-adjustment ledger
through a canonical nine-activity loan scenario. The loan has an initial
principal of \$1,000,000 with an annual interest rate of 10.0\%,
computed on an actual/365 day-count basis (daily rate: \$273.97).

\begin{table}[h!]
\centering
\caption{Nine-activity evaluation scenario.}
\label{tab:scenario}
\begin{tabular}{clrl}
  \toprule
  \textbf{Step} & \textbf{Activity} & \textbf{Amount} & \textbf{Effective Date} \\
  \midrule
  1 & Disbursal              & \$1,000,000 & 2024-03-01 \\
  2 & SOD (interest accrual) & ---         & 2024-03-02 \\
  3 & SOD (interest accrual) & ---         & 2024-03-03 \\
  4 & Payment \#234          & \$10,000    & 2024-03-03 \\
  5 & SOD (interest accrual) & ---         & 2024-03-04 \\
  6 & SOD (interest accrual) & ---         & 2024-03-05 \\
  7 & Payment \#456          & \$5,000     & 2024-03-05 \\
  8 & Reversal of \#234      & ---         & 2024-03-05 \\
  9 & Backdated Payment \#768 & \$4,000   & 2024-03-03 (recv.\ 2024-03-05) \\
  \bottomrule
\end{tabular}
\end{table}

Steps 1--7 are in-order activities processed identically by both
systems. Steps 8 and~9 trigger the retroactive paths.

\paragraph{Conventional ledger baseline.}
Our comparative baseline follows typical loan servicing practice:
reversals are handled by appending a compensation entry plus a single
grouped downstream adjustment; backdated payments are appended against
current processing-time balances. It correctly maintains running
balances but represents multiple downstream corrections as a single
entry.

\subsection{Comparison at Activity 8 (Reversal)}

After processing steps 1--7, the primary ledger for both systems holds
7~entries with identical balances: principal \$990,547.24; interest
\$0.00 (fully consumed by Payment~\#456).

\paragraph{Conventional ledger (step~8).}
Appends a compensation entry ($+$\$9,452.05 principal,
$+$\$547.95 interest) followed by a \emph{single grouped adjustment}
($+$\$3.30 principal, \$0.00 net interest) covering three distinct
downstream corrections: (i)~SOD-304 interest recalculation,
(ii)~SOD-305 interest recalculation, and (iii)~Payment~\#456
re-allocation. Total entries after reversal: \textbf{9}.

\paragraph{Fork-and-sync ledger (step~8).}
Applies the same compensation entry. The fork-replay exposes that
SOD-304 and SOD-305 each recalculated on a higher principal (restored
to \$1,000,000) and Payment~\#456 required re-allocation with a higher
interest base. Three individually traced adjustment entries are
appended:
\begin{enumerate}
  \item Adjustment for SOD-20240304: $\Delta\text{interest} = +\$1.65$
        (interest: \$271.38 $\to$ \$273.03)
  \item Adjustment for SOD-20240305: $\Delta\text{interest} = +\$1.65$
        (same recalculation)
  \item Adjustment for Transaction-456: $\Delta\text{principal} = +\$3.30$,
        $\Delta\text{interest} = -\$3.30$ (re-allocation to higher
        interest base)
\end{enumerate}
Total entries after reversal: \textbf{11}. Every correction row names
the entry it corrects; the ledger is readable top-to-bottom with no
opaque rows.

\subsection{Comparison at Activity 9 (Backdated Payment)}

\paragraph{Conventional ledger (step~9).}
The \$4,000 payment is applied against the current processing-time
interest balance of \$547.95. Allocation: \$547.95 to interest,
\$3,452.05 to principal (one payment entry). The revised principal
lowers the basis for Payment~\#456's prior allocation, so the system
appends one additional grouped re-allocation adjustment entry.
Total: 9 + 2 = \textbf{11} entries.

\paragraph{Fork-and-sync ledger (step~9).}
The payment is inserted at its economic date of 2024-03-03. At that
date, the interest balance (per the corrected retroactive fork) is
\$547.94. Allocation: \$547.94 to interest, \$3,452.06 to principal.
The \$0.01 difference is not drift; it reflects the economically
correct interest basis at the payment's effective date (R1). After
replay, downstream SODs (March~4 and~5) recalculate on the lower
post-payment principal. The delta per SOD is less than \$0.01 per day
when rounded, so those two SOD adjustment entries are suppressed
(zero-delta suppression). Payment~\#456's re-allocation, however,
changes non-trivially against the revised interest basis, generating
one additional adjustment entry. Total: 11 + 2 = \textbf{13} entries.

\subsection{Summary Results}

\begin{table}[h!]
\centering
\caption{Comparison summary after all nine activities.}
\label{tab:summary}
\begin{tabularx}{\linewidth}{Xcc}
  \toprule
  \textbf{Metric} & \textbf{Conventional} & \textbf{Fork-and-Sync} \\
  \midrule
  Final principal balance     & \$987,095.18 & \$987,095.19 \\
  Final interest balance      & \$0.00       & \$0.00 \\
  Total entries (all 9 steps) & 11           & 13 \\
  Downstream correction entries (steps 8\,\&\,9) & 2 grouped  & 4 individually traced \\
  Entry-level traceability    & No           & Yes \\
  Retroactive composability   & Fragile      & Yes \\
  Grouped-entry decomposition & Not possible & Not applicable \\
  \bottomrule
\end{tabularx}
\end{table}

The fork-and-sync approach generates two additional entries (11 vs.\ 13)
as the cost of full traceability. This is $O(k)$ entry overhead per
retroactive event---proportional to the number of divergent downstream
activities~$k$, not to the total ledger length~$n$. In practice $k$ is
a small constant (Section~\ref{sec:discussion}), making the overhead
effectively bounded and predictable.

\subsection{Retroactive Composability}

Consider a hypothetical tenth activity: a second backdated event at
2024-03-04. With grouped corrections, the conventional engine must
reason about adjusting-an-adjustment: the grouped $+\$3.30$ entry from
step~8 becomes the target of further correction, and its decomposition
into three original effects (SOD-304, SOD-305, Payment~\#456) cannot be
recovered from the ledger alone. Full re-derivation from the primary
activity log is required.

With per-entry traced corrections, each prior adjustment row identifies
its exact source activity. The fork at 2024-03-04 simply re-encounters
these source activities during replay, and the
\texttt{calculateTotalImpact} / \textsc{Sync} protocol naturally
converges to the new correct state. No special-case decomposition logic
is needed.

% =========================================================================
\section{Discussion}
\label{sec:discussion}
% =========================================================================

\subsection{Entry Count Growth}

The fork-and-sync approach generates strictly more ledger entries than
the grouped-adjustment approach---$k$~individual adjustment rows versus
1 grouped row, where $k$ is the number of downstream activities with
divergent impacts. In practice, $k$ is small (typically 3--10 for a
standard loan ledger), so growth is bounded and predictable.

\subsection{Zero-Delta Suppression}

The skip-if-equal condition in Algorithm~\ref{alg:sync} (Step~15) is
important for practical usability. Without it, every activity after the
fork point would produce an adjustment entry, inflating the ledger by
$O(n')$ entries per retroactive event (where $n'$ is the total number of
replay activities after the fork point). With suppression, only the $k$
activities with genuine downstream impact (typically those that use the
running balance in their computation, such as temporal accruals) produce
entries. Allocation-fixed activities (disbursals, static-fee charges)
almost never diverge after a principal change.

\subsection{Virtual Clock and Determinism}

The \texttt{LedgerClock} design is a form of test-seam
injection~\cite{feathers2004} applied at the production-logic level
rather than only in tests. By parameterizing all time-sensitive
computations on an external clock object, the system achieves
reproducibility: replaying the same activity sequence with the same
clock configuration always produces the same ledger. This is a
prerequisite for regression testing of retroactive paths, which are
notoriously difficult to test with wall-clock-coupled logic.

\subsection{Limitations}

\begin{itemize}
  \item \textbf{Concurrent retroactive events.}~The current design
        assumes serial application of activities to a given ledger.
        Concurrent retroactive events targeting the same ledger require
        external serialization (e.g., optimistic locking on the ledger
        object) to prevent lost-update anomalies~\cite{grayreuter1992}.
  \item \textbf{Infinite recursion guard.}~The reversal algorithm
        excludes the target activity from replay explicitly. Incorrect
        activity-ID matching or missing exclusion in a derived
        implementation would cause unbounded recursion. The system
        mitigates this via exhaustive fixture-based tests.
  \item \textbf{Query performance.}~\texttt{calculateTotalImpact}
        performs a linear scan over ledger entries filtered by source
        activity. For long ledgers, a secondary composite index on
        (\texttt{loanId}, \texttt{sourceLedgerActivityType},
        \texttt{sourceLedgerActivityId}) is required to maintain
        $O(\log n)$ query cost~\cite{graefe2011}.
  \item \textbf{Single-entry model.}~The current model does not enforce
        double-entry bookkeeping constraints (debits~=~credits).
        Extension to double-entry contra-account ledgers is future work.
\end{itemize}

\subsection{Future Work}

Future directions include: (i)~a formal proof of convergence and
correctness for recursively nested retroactive events;
(ii)~benchmarking on synthetic ledgers of 10,000--100,000 entries to
validate the $O(p + m \log m)$ empirical cost model;
(iii)~extension to multi-party ledgers where a retroactive event on one
party's ledger triggers cascading reconciliation to a counterparty
ledger.

% =========================================================================
\section{Conclusion}
\label{sec:conclusion}
% =========================================================================

We presented a fork-and-sync architecture for financial ledger systems
that handles both out-of-order (backdated) activities and reversals
while satisfying all four requirements: correctness, audit immutability,
entry-level traceability, and retroactive composability.
The key contribution is Algorithm~\ref{alg:sync} (\textsc{Sync}), which
reconciles per-activity balance impacts between a retroactive fork and
the primary ledger, emitting one individually source-attributed
adjustment entry per divergent activity rather than a single grouped
correction. A virtual ledger clock ensures deterministic replay of
time-dependent derived calculations. Evaluation on a nine-activity loan
scenario demonstrates that the approach produces an auditable,
row-by-row explainable ledger with a bounded entry-count overhead of
$O(k)$ per retroactive event, where $k$ is the number of activities
with genuine downstream balance impact. A complete reference
implementation is publicly available. We believe this approach is
directly applicable to any financial system that must accommodate
retroactive corrections over a sequence of derived calculations,
including general ledger systems, derivative position accounts, and
insurance claim ledgers.

% =========================================================================
\section*{Acknowledgements}
% =========================================================================

The authors thank the open-source contributors to the Quarkus,
Hibernate, and JavaMoney/Moneta projects, whose frameworks underpin the
reference implementation.

% =========================================================================
\begin{thebibliography}{99}

\bibitem{fowler2005}
M.~Fowler, ``Event Sourcing,'' 2005.
\url{https://martinfowler.com/eaaDev/EventSourcing.html}.

\bibitem{snodgrass1999}
R.~T.~Snodgrass,
\textit{Developing Time-Oriented Database Applications in SQL}.
Morgan Kaufmann, 1999.

\bibitem{jensen1994}
C.~S.~Jensen, J.~Clifford, R.~Elmasri, S.~K.~Gadia, P.~Hayes,
and S.~Jajodia,
``A Glossary of Temporal Database Concepts,''
\textit{ACM SIGMOD Record}, vol.~23, no.~1, pp.~52--64, 1994.

\bibitem{garciamolinasalem1987}
H.~Garcia-Molina and K.~Salem,
``Sagas,''
in \textit{Proc. ACM SIGMOD International Conference on Management of Data},
pp.~249--259, 1987.

\bibitem{bernsteingoodman1981}
P.~A.~Bernstein and N.~Goodman,
``Concurrency Control in Distributed Database Systems,''
\textit{ACM Computing Surveys}, vol.~13, no.~2, pp.~185--221, 1981.

\bibitem{grayreuter1992}
J.~Gray and A.~Reuter,
\textit{Transaction Processing: Concepts and Techniques}.
Morgan Kaufmann, 1992.

\bibitem{feathers2004}
M.~C.~Feathers,
\textit{Working Effectively with Legacy Code}.
Prentice Hall, 2004.

\bibitem{graefe2011}
G.~Graefe,
``Modern B-Tree Techniques,''
\textit{Foundations and Trends in Databases}, vol.~3, no.~4,
pp.~203--402, 2011.

\end{thebibliography}

% =========================================================================
\appendix
% =========================================================================

\section{Test Fixture Summary}

The reference implementation includes fixture-based integration tests
that encode final ledger expectations as CSV files. The primary fixture
for the nine-activity evaluation scenario is:

\begin{lstlisting}
src/test/resources/data/ledger/service/
  applyLedgerActivities_test1/
    apply_ledger_activities_test1_expectation.csv
\end{lstlisting}

Each row encodes a full ledger entry expectation including entry type,
source activity ID, incremental and running balances for all four
components. These fixtures serve as executable specifications equivalent
to the comparison in Table~\ref{tab:summary} and are verified on every
CI run.

\section{ACM Computing Classification System (CCS 2012)}
\label{app:ccs}

\begin{itemize}
  \item \textbf{Primary:} Information systems $\to$ Transaction management
  \item \textbf{Secondary:} Information systems $\to$ Data streams;
        Software and its engineering $\to$ Software reliability
\end{itemize}

\noindent\textit{CCS format:}
\texttt{Information systems\textasciitilde Transaction management;
Information systems\textasciitilde Data streams;
Software and its engineering\textasciitilde Software reliability}

\end{document}
